\chapter{Validation and Results}
\label{chap:validation}

Chapters 5 and 6 described what the platform is built to do and how it is
built. This chapter asks a different question, whether it actually does what
it claims, and how this was checked. I present the three layers of testing
that support the deterministic engine and the platform built around it, walk
through one concrete assessment profile from raw score to final score,
evaluate the local AI layer against its own guardrails, and close with the
limitations I consider important to state explicitly rather than leave for a
reader to discover on their own.

\section{Validation strategy}
\label{sec:validation-strategy}

The validation approach follows directly from the inviolable principles
introduced in Chapter 4 (\ref{chap:method}), in particular the requirement
that the scoring engine be deterministic and auditable. Testing is organised
in three layers, moving from the pure computation core outward to the full
request path a user actually exercises.

\subsection{Three layers of testing}
\label{subsec:three-layers}

The first layer tests the scoring engine in isolation, as a pure Python
package with no I/O. The second layer tests the FastAPI backend that wraps
it, exercising HTTP endpoints against a real (in-memory) database. The third
layer is a targeted cross-check of the single seam between the two, the
translation from stored, raw answers to the engine's input format.

\begin{table}[htbp]
  \centering
  \begin{tabular}{@{}lll@{}}
    \toprule
    Layer & Scope & Location \\
    \midrule
    Engine unit tests & Pure scoring pipeline, no I/O & \texttt{app/scoring/tests/} \\
    Backend integration tests & FastAPI endpoints, async DB & \texttt{backend/tests/} \\
    End-to-end bridge cross-check & DB row to \texttt{AssessmentInput} to result & \texttt{score\_test.py} \\
    \bottomrule
  \end{tabular}
  \caption{The three testing layers}
  \label{tab:testing-layers}
\end{table}

\begin{figure}[htbp]
  \centering
  \begin{tikzpicture}[
      >=Latex,
      node distance=4mm,
      proc/.style={
        draw=black!65, thick, rounded corners=2pt, fill=black!4,
        minimum height=1.0cm, text width=1.9cm, align=center,
        font=\footnotesize
      },
      data/.style={
        draw=black!55, thick, dashed, rounded corners=2pt, fill=white,
        minimum height=1.0cm, text width=1.9cm, align=center,
        font=\footnotesize\itshape
      },
      band/.style={
        rounded corners=2pt, thick, minimum height=0.85cm
      },
      guide/.style={draw=black!35, dashed, line width=0.5pt},
      lbl/.style={font=\scriptsize, align=center}
    ]

    \node[proc]                          (n1) {HTTP\\request};
    \node[proc, right=of n1]             (n2) {FastAPI\\endpoint + DB};
    \node[proc, right=of n2]             (n3) {scoring\_\\bridge.py};
    \node[data, right=of n3]             (n4) {Assessment\\Input};
    \node[proc, right=of n4]             (n5) {Engine\\(Ch.~5)};
    \node[data, right=of n5]             (n6) {Assessment\\Result};
    \node[proc, right=of n6]             (n7) {HTTP\\response};

    \draw[->] (n1) -- (n2);
    \draw[->] (n2) -- (n3);
    \draw[->] (n3) -- (n4);
    \draw[->] (n4) -- (n5);
    \draw[->] (n5) -- (n6);
    \draw[->] (n6) -- (n7);

    \draw[guide] (n4.south west) -- ++(0,-0.9);
    \draw[guide] (n6.south east) -- ++(0,-0.9);
    \draw[band, draw=blue!55!black, fill=blue!8]
      ($(n4.south west)+(0,-0.9)$) rectangle ($(n6.south east)+(0,-1.75)$);
    \node[lbl, blue!55!black] at ($(n4.south east)!0.5!(n6.south west)+(0,-1.325)$)
      {Engine unit tests};

    \draw[guide] (n2.south west) -- ++(0,-2.0);
    \draw[guide] (n4.south east) -- ++(0,-2.0);
    \draw[band, draw=teal!60!black, fill=teal!8]
      ($(n2.south west)+(0,-2.0)$) rectangle ($(n4.south east)+(0,-2.85)$);
    \node[lbl, teal!60!black] at ($(n2.south east)!0.5!(n4.south west)+(0,-2.425)$)
      {End-to-end bridge\\cross-check};

    \draw[guide] (n1.south west) -- ++(0,-3.1);
    \draw[guide] (n7.south east) -- ++(0,-3.1);
    \draw[band, draw=black!55, fill=black!5]
      ($(n1.south west)+(0,-3.1)$) rectangle ($(n7.south east)+(0,-3.95)$);
    \node[lbl, black!55] at ($(n1.south east)!0.5!(n7.south west)+(0,-3.525)$)
      {Backend integration tests};

  \end{tikzpicture}
  \caption{Overview of the validation pipeline}
  \label{fig:validation-pipeline}
\end{figure}

The engine suite grew from 64 tests at the point of the independent audit
(\ref{chap:scoring-engine}) to 67 after the corrections it triggered, three
new tests specifically targeting a gap the audit had identified (see
\ref{subsec:test-profiles}), and stands at 70 today. The backend suite,
which exercised only the two endpoints required by the end of Sprint 0,
\texttt{test\_health.py} and \texttt{test\_referential\_endpoint.py}, grew
alongside every backend feature landed afterward, assessments, scoring,
authentication and access control, and now counts 68 tests.

\subsection{Deterministic testing philosophy}
\label{subsec:deterministic-philosophy}

Because the engine is deterministic by construction (principle P1, Chapter
4), its tests can assert exact values rather than tolerance ranges. A test
such as \texttt{test\_determinism}, which runs the same assessment through
the engine twice and compares the two serialized outputs, would be
meaningless for a component with any randomness or hidden state. This is
also why the golden scenarios described below assert precise numeric scores
rather than approximate bounds, and why the boundary sweep test can claim
exhaustive coverage of a parameter range rather than sampling it. This
certainty does not extend to the AI layer, which is evaluated on different
terms in \ref{sec:ai-evaluation}.

\subsection{Test profiles built for coverage}
\label{subsec:test-profiles}

Beyond the determinism check, the engine suite is organised around a small
number of named scenarios chosen to exercise specific parts of the scoring
pipeline described in Chapter 5, particularly the parts most likely to hide
a subtle error, ordering, thresholds, and edge conditions.

\begin{table}[htbp]
  \centering
  \begin{tabular}{@{}lp{9cm}@{}}
    \toprule
    Test & Purpose \\
    \midrule
    \texttt{test\_determinism} & Same input run twice yields identical output \\
    Golden scenarios (5) & End-to-end reference cases covering the full pipeline \\
    \texttt{test\_bounds\_across\_sweep} & 20 parametrized combinations, proves the final score stays in $[0, 100]$ \\
    \texttt{test\_evidence\_caps\_are\_strictly\_increasing} & Evidence-type caps are monotonic \\
    \texttt{test\_golden\_na\_is\_neutral} & A N/A best practice contributes to neither numerator nor denominator \\
    Loader validation tests & Missing file, malformed YAML, missing field all raise \\
    \texttt{test\_golden\_consistency\_penalty\_triggered} & A maturity gap above threshold deducts the expected points \\
    \texttt{test\_golden\_two\_gates\_strictest\_cap\_wins} & Two simultaneous foundational gates, the stricter cap is retained \\
    \texttt{test\_penalty\_threshold\_is\_strict\_greater\_than} & A gap exactly at the threshold does not penalize, just above does \\
    \bottomrule
  \end{tabular}
  \caption{Named test scenarios and what each one locks down}
  \label{tab:test-profiles}
\end{table}

The last three rows were added after an independent audit of the engine
found that the consistency-penalty branch, one of the two original
correction mechanisms of the model, was never exercised by any test (finding
H1, \ref{chap:scoring-engine}). This is worth dwelling on briefly. A
64-test suite is a good headline number, but it had a gap precisely where
the model's own differentiating logic lived. The three tests added close
that gap without touching a single line of production code, they capture
the pipeline's existing, correct behaviour rather than changing it.

\section{Functional results}
\label{sec:functional-results}

Where Section~\ref{sec:validation-strategy} described how the platform is
tested, this section shows what those tests actually establish, first about
the seam between storage and computation, then about the scoring pipeline
itself through a single worked example.

\subsection{API-engine cross-check}
\label{subsec:bridge-crosscheck}

The engine suite (Chapter~\ref{chap:scoring-engine}) proves that the scoring
pipeline is correct in isolation, given a well-formed \texttt{AssessmentInput}.
It says nothing about whether the platform actually builds a well-formed
\texttt{AssessmentInput} from what a user submitted. That translation is the
single most likely place for a silent mismatch to appear, a renamed field, a
dropped answer, a best practice marked N/A on one side and not the other, so
it is deliberately isolated in one module, \texttt{scoring\_bridge.py}, rather
than spread across the API layer.

\texttt{scoring\_bridge.py} is the only module in the codebase that imports
both the SQLAlchemy models and the engine's Pydantic input models. Its
\texttt{build\_assessment\_input} function maps a stored \texttt{Assessment}
row, with its answers and best-practice annotations eagerly loaded, into an
\texttt{AssessmentInput}, and \texttt{missing\_answers} independently
precomputes the full list of unanswered questions.

This second function exists because of a risk raised during the engine audit
(Chapter~\ref{chap:scoring-engine}), before the bridge was written, an
incomplete assessment would reach the engine, which raises a bare
\texttt{KeyError} on the first missing answer it encounters rather than
reporting the assessment as incomplete. Left unhandled, this would have
surfaced as an opaque \texttt{500} response. \texttt{missing\_answers} closes
that gap ahead of time, so the scoring endpoint can return every missing
answer at once as a \texttt{422}, keeping genuine user errors out of the
server-error range, as documented in \ref{chap:platform}. This is a small
illustration of a pattern that recurs through the project, a risk identified
on paper is treated as a requirement on the next component, not as an
afterthought once it fails in practice.

\texttt{score\_test.py} exercises this seam end to end, it builds an
assessment through the API, submits answers, requests a score, and checks the
result against what calling the engine directly on the same input would
produce, proving that the bridge introduces no distortion.

\subsection{Worked example: from raw score to final score}
\label{subsec:worked-example}

The scoring pipeline applies, in order, an evidence cap, a phase-weighted
aggregation, a foundational-gate cap, and a consistency-penalty deduction
(Chapter~\ref{chap:scoring-engine}). Each step is unit-tested independently
(Section~\ref{subsec:test-profiles}), but a single assessment profile that
exercises several of them in sequence makes the cumulative effect concrete in
a way that isolated tests do not. This is also, deliberately, the numeric
proof that the engine audit found missing (finding H1,
Chapter~\ref{chap:scoring-engine}), the correction added the test, this
section reports what it actually shows.

Table~\ref{tab:worked-example} follows the \texttt{penalized} test profile
through the pipeline. The four phase scores it produces, 75 for Prevention,
68 for Preparedness, 75 for Response, and 36 for Recovery, average to a raw
global score of 63 under the referential's equal 25 percent phase weighting.
BP08, a foundational best practice, scores 0 on this profile, well below its
40-point threshold, which triggers its foundational gate and caps the global
score at 65. Because the raw score of 63 already sits below that cap, the
gate is recorded as applied but does not move the number, an instructive
edge case, a gate can fire without being the binding constraint. Three
consistency pairs then trigger simultaneously, BP10 against BP11, BP07
against BP08, and BP05 against BP15, each with a maturity gap of 4.0 against
an allowed maximum of 2, each deducting 5 points. The three penalties bring
the score down by 15 points in total, from 63 to a final global score of 48
(48.4 before rounding), mapped to maturity level 2, Defined.

\begin{table}[htbp]
  \centering
  \begin{tabular}{@{}lr@{}}
    \toprule
    Stage & Score \\
    \midrule
    Raw global score & 63 \\
    After foundational gate & 63 \\
    After consistency penalties & 48 \\
    Final global score & 48 \\
    \bottomrule
  \end{tabular}
  \caption{Worked example: the \texttt{penalized} profile}
  \label{tab:worked-example}
\end{table}

This example is a useful counterpoint to a purely additive reading of the
model. A respondent looking only at individual best-practice scores could
reasonably expect a global score closer to the mid-sixties. The gap between
that expectation and the actual 48 is exactly the point of the two
correction mechanisms, a single very weak foundational best practice and a
pattern of inconsistent maturity across related practices both leave a
visible mark on the final number, which is the intended, auditable
behaviour rather than a computational surprise.

\section{AI evaluation}
\label{sec:ai-evaluation}

The scoring engine can be validated the way Sections~\ref{sec:validation-strategy}
and~\ref{sec:functional-results} did it, against exact expected values, because
it is deterministic. The local AI layer is not, and evaluating it calls for
different terms, the reasoning behind the model choice, the guardrails put in
place because the model cannot be trusted unsupervised, how often those
guardrails have to intervene, and a qualitative read of what the output is
actually worth.

\subsection{Model selection: qwen2.5:3b versus a 7B model}
\label{subsec:model-selection}

Ollama is configured to run \texttt{qwen2.5:3b} by default rather than a
larger model such as \texttt{mistral:7b-instruct}, one of the candidates
originally listed for this role (\texttt{referential.yaml},
\texttt{ai\_module\_scope}). The reasoning is documented in the architecture
log, inference runs on CPU inside the demo VM, where a 7B model's latency
made the recommendations endpoint impractical, multi-minute responses
observed informally rather than through a controlled benchmark, while
\texttt{qwen2.5:3b} keeps generation to approximately one minute. This is
defensible because the task given to the model is constrained rather than
open-ended, it summarizes and enriches recommendations that the referential
already supplies for a handful of weak best practices, it does not reason
freely over the whole framework.

Even at roughly a minute, blocking the dashboard on generation was not
acceptable, which is why recommendations are fetched asynchronously and
cached in memory per assessment, the dashboard renders immediately and the
narrative fills in once ready, fetched once and reused across navigation
rather than regenerated on every visit (\ref{chap:platform}). The model
choice and this caching decision address the same constraint from two
directions, one lowers the cost of a generation, the other lowers how often a
generation has to happen at all.

The trade-off is left open rather than closed, the model is swapped through a
single environment variable, and a controlled, timed comparison against a
larger model remains future work (\ref{subsec:perspectives-future-work}).

\subsection{Failure modes and guardrails}
\label{subsec:failure-modes}

A 3B model prompted to reference specific best-practice identifiers is not a
reliable source of exact strings, so the AI layer is built on the assumption
that it will occasionally get an identifier wrong rather than on the hope
that it will not. The scoring role is fixed regardless, the model receives a
fully computed \texttt{AssessmentResult} and produces text only, it never
touches a score (principle P1, \ref{chap:scoring-engine}). The prompt is
anchored to the exact list of best-practice identifiers relevant to the
assessment and carries no personally identifiable information, no
organization name, no identifiers, no N/A justification text (principle P3).

Two guardrails sit downstream of generation, before any text reaches the
user. First, identifier normalization corrects the shapes a small model
tends to produce, a bare \texttt{BP4} is padded to \texttt{BP04}, but a fused
identifier such as \texttt{BP123} is rejected outright rather than guessed
apart into \texttt{BP01} and \texttt{BP23}, since splitting it would be
fabricating a claim the model never actually made about either practice.
Second, every normalized identifier is checked for existence against the
loaded referential, an action citing a best practice that is not in the
framework, such as a fabricated \texttt{BP18} on a referential that only
defines sixteen, is discarded rather than shown.

The result is that a hallucinated identifier or a malformed response is a
degraded response, not an incorrect one, the worst case is losing the
AI-generated narrative for that best practice, never seeing a fabricated one.
This existence check is scoped to the full referential rather than to the
narrower subset of weak best practices sent in a given prompt, a distinction
worth naming explicitly in Section~\ref{subsec:ai-quality-limitation}.

\subsection{Residual fallback rate}
\label{subsec:fallback-rate}

The guardrails in Section~\ref{subsec:failure-modes} are only as reassuring
as they are infrequently needed, a system that falls back on every request
would be technically safe but practically equivalent to not having an AI
layer at all. This frequency was not tracked systematically during
development, no instrumentation logged how often the endpoint actually
returned \texttt{source="fallback"} across real or test invocations, so no
rate can be reported here. The fallback path was directly observed at least
once during testing, for the assessment discussed in
Section~\ref{subsec:worked-example}, which returned the referential's
static recommendations rather than an AI-generated narrative
(Section~\ref{subsec:qualitative-ai}). One observed instance says nothing
about frequency, but it does confirm the fallback path is exercised in
practice and not merely a theoretical branch. This absence of a tracked rate
is itself worth naming rather than papering over, and it points to concrete
follow-up work, instrumenting the recommendations endpoint to log its own
degradation rate would turn this qualitative observation into a measurable
one. Once a network misconfiguration between the demo VM and the host
running Ollama was diagnosed and corrected, a successful AI-generated
response was also obtained on the same profile, confirming that both the
fallback path and the AI path are exercised in practice rather than one
being purely theoretical (Section~\ref{subsec:qualitative-ai}).

\subsection{Qualitative assessment of recommendations}
\label{subsec:qualitative-ai}

The referential's static text is a fixed sentence per best practice,
independent of how the organization actually answered
(\texttt{fallback\_recommendation\_when\_low\_maturity},
\texttt{referential.yaml}). The AI layer's stated responsibility is to go
beyond that, to produce an executive synthesis in natural language, to
prioritize the three to five actions that matter most across the weak best
practices found, taking foundational status and dependencies into account,
and to contextualize each recommendation using what the user actually
answered rather than the generic text alone
(\texttt{ai\_module\_scope.downstream\_responsibilities},
\texttt{referential.yaml}).

Table~\ref{tab:recommendation-comparison} compares both outputs on the same
\texttt{penalized} profile discussed throughout this chapter, static referential
text against the AI-generated narrative for the same three priority best
practices, BP08, BP11, and BP15.

\begin{table}[htbp]
  \centering
  \small
  \begin{tabular}{@{}p{1.8cm}p{7cm}p{7cm}@{}}
    \toprule
    Field & Static (referential fallback) & AI-generated (\texttt{qwen2.5:3b}) \\
    \midrule
    Executive summary &
    Overall crisis-management maturity is assessed at 48.4/100, level 2,
    Defined. Focus first on the priority actions below, which target the
    weakest and most structural best practices. &
    The organization's cyber crisis management is below par, with scores
    averaging at level 2, Defined. Key priorities include establishing
    clear escalation criteria for timely activation of plans, implementing
    a comprehensive training program for cyber crisis managers, and
    developing robust Business Resumption Plans. \\
    \addlinespace
    BP08 &
    Define operational escalation criteria, severity, scope, time
    pressure, translate into decisions, who, when, how, and test in
    exercises. Foundational best practice, score 0.0/100. &
    Define operational escalation criteria including severity, scope, and
    time pressure, and test in exercises. Timely activation is crucial to
    effective response, establishing clear guidelines will improve
    readiness. \\
    \addlinespace
    BP11 &
    Implement a training curriculum for cyber crisis roles, including
    onboarding, regular sessions, and effectiveness measurement. Score
    0.0/100. &
    Implement a training curriculum for cyber crisis roles with regular
    sessions and effectiveness measurement. Continuous training ensures
    staff are prepared, aligning responsibilities with escalation
    criteria. \\
    \addlinespace
    BP15 &
    Develop an integrated BRP with BCP/DR, regular review, recovery tests,
    and stakeholder involvement beyond IT. Score 0.0/100. &
    Develop an integrated Business Resumption Plan aligned with the
    Business Continuity Plan and Disaster Recovery Plan, including regular
    reviews and stakeholder involvement beyond IT. Resilience planning is
    critical for recovery, a comprehensive BRP ensures a coordinated
    approach to business continuity. \\
    \bottomrule
  \end{tabular}
  \caption{Static referential text against AI-generated narrative, same
  profile, same three best practices}
  \label{tab:recommendation-comparison}
\end{table}

The comparison is more modest than the stated ambition. For all three
actions, the AI output is close to a paraphrase of the static text rather
than a substantively new recommendation, BP08 keeps the same three criteria
in the same order, BP11 keeps the same structure with one clause dropped,
BP15 mostly expands the acronyms. What the AI genuinely adds is a short
justification sentence appended to each action, absent from the static
version, explaining in general terms why the action matters. This is a real
addition, but a generic one, the justification for BP08 or BP11 would read
just as naturally attached to a different weak best practice, it is not
visibly grounded in this organization's specific pattern of failure, three
triggered consistency penalties on top of one marginal foundational gate
(\ref{subsec:worked-example}).

The executive summary shows a comparable trade, the AI version drops the
precise 48.4 score present in the static text in favour of a looser
paraphrase, "scores averaging at level 2", which is imprecise in a way the
static text is not, the final score is not a simple average, it results from
the gate and penalty mechanism described in
Section~\ref{subsec:worked-example}. No identifier was fabricated and no
best practice outside the referential's sixteen was cited in this instance,
so the guardrails held, but holding the guardrails and adding genuine
analytical value are two different bars, and on this single example the
output clears the first more convincingly than the second.

\section{Limitations}
\label{sec:limitations}

The guardrails and cross-checks described in the previous sections make the
platform trustworthy within a defined scope. They do not make the underlying
model of maturity itself complete. Four limitations are worth stating plainly
rather than leaving for a reader to find, because each one bears directly on
how much weight a result from this tool should be given.

\subsection{Self-assessment is inherently declarative}
\label{subsec:declarative-limitation}

Every score the platform produces ultimately rests on what a respondent
chooses to answer. Nothing in the data model requires, or is able to require,
that a declared maturity level of 3 for a given question actually reflects
the organization's practice rather than the respondent's understanding of it,
or their incentive to present it favourably. This is not a flaw specific to
this implementation, it is a property of self-assessment as a method, shared
with essentially every maturity framework that does not involve an external
auditor collecting evidence directly. The platform inherits it rather than
solves it, and the two mitigations built against it, evidence caps and the
consistency penalties, are best understood as bounding the damage a purely
declarative answer can do to the score, not as removing the underlying
dependency on honest, informed self-reporting.

\subsection{The evidence cap mitigates declarativeness only partially}
\label{subsec:evidence-cap-limitation}

The evidence-type mechanism (\ref{chap:scoring-engine}) caps a declared
maturity level according to the strongest evidence category the respondent
selects, an unsupported claim cannot score above level 1, a claim backed by a
real incident can reach level 4. This meaningfully narrows how far a purely
declarative answer can inflate a score. It does not close the gap entirely,
because the evidence type itself is self-selected in exactly the same way
the maturity level is, the data model records which category the respondent
chose, not any independent proof that the category is accurate. A respondent
willing to overstate a maturity level is equally free to overstate the
evidence behind it, and the platform has no upload, audit trail, or
verification step that would catch this. The cap changes the shape of the
incentive to inflate a score, it does not remove it.

\subsection{No inter-rater reliability study}
\label{subsec:interrater-limitation}

The scoring engine's determinism guarantees that the same answers always
produce the same score (\ref{chap:scoring-engine}). It says nothing about
whether two different, equally informed people answering on behalf of the
same organization would give the same answers in the first place. No study
was run to check this, no assessment was completed independently by more
than one respondent from the same organization and compared. Without that
data, a claim that the tool measures maturity objectively, rather than
measuring one respondent's perception of it consistently, would be
overstated. This is a natural extension for validating the methodology
further, not a defect in what was built during the internship.

\subsection{AI output quality not formally measured beyond anchoring constraints}
\label{subsec:ai-quality-limitation}

Section~\ref{sec:ai-evaluation} established that the AI layer cannot corrupt
a score and cannot cite a best practice outside the referential, both
verified properties. That second guarantee is narrower than it first
appears, it checks membership in the sixteen-best-practice referential as a
whole, not membership in the smaller set of weak best practices actually
sent to the model for a given assessment (\ref{subsec:failure-modes}). A
real but contextually irrelevant identifier would currently pass this check
undetected, a gap that stops short of a hallucinated global score or a
fabricated best practice, but is narrower than the guarantee as informally
described.

Neither guarantee says anything about whether the generated narrative is
actually good advice, whether the prioritization across weak best practices
is sound, whether the contextualization beyond the static fallback text is
genuinely more useful to a reader than the fallback would have been on its
own. The single worked comparison in Section~\ref{subsec:qualitative-ai}
suggests the current output leans closer to paraphrase than to genuine
contextualization, but one example is an illustration, not a formal
evaluation. Evaluating that properly would require either a domain expert
scoring a larger sample of generated recommendations against some rubric, or
feedback from someone who used the tool in a real advisory context, neither
of which was carried out. What Section~\ref{sec:ai-evaluation} demonstrates
is that the AI layer is safe to expose, not that its output has been shown
to be good, and the two claims should not be conflated.

Taken together, these four limitations describe the boundary of what this
internship validated, and what it did not have the time or the setting to
validate. None of them undermine the platform's core claim, that a
deterministic, auditable score can be computed from structured self-reported
answers, but they do bound how that score should be interpreted and by whom.
Chapter~\ref{chap:pm} returns to several of these points from a project
management angle, in particular what a next iteration would prioritize given
more time.